\documentclass[12pt]{article}
\usepackage{amsmath,amssymb,amsfonts}
\usepackage[margin=1in]{geometry}

\begin{document}

\section*{Chapter 2, Problem 5}

\subsection*{Problem Statement}
Demonstrate the following reciprocity relationship for the simple isoperimetric problem: The particular function which renders
\[
  I = \int_{x_1}^{x_2} f(x, y, y')\,dx
\]
an extremum with respect to functions which give
\[
  J = \int_{x_1}^{x_2} g(x, y, y')\,dx
\]
a prescribed value, also renders $J$ an extremum with respect to functions which give $I$ a prescribed value if the Lagrange multiplier $\lambda \neq 0$.

In all our discussions so far on finding the function $f$ for which $I = \int f\,dx$ is an extremum, it has been assumed that $f$ depends only on $x$, $y$, and $y'$.

\subsection*{Solution}

In the isoperimetric problem, we seek to extremize $I$ subject to the constraint $J = \text{const}$. We introduce a Lagrange multiplier $\lambda$ and form the augmented functional:
\[
  K = I + \lambda J = \int_{x_1}^{x_2} [f + \lambda\,g]\,dx.
\]
The necessary condition for an extremum is the Euler--Lagrange equation for $h = f + \lambda g$:
\[
  \frac{\partial h}{\partial y} - \frac{d}{dx}\frac{\partial h}{\partial y'} = 0,
\]
which is:
\begin{equation}
  \frac{\partial f}{\partial y} - \frac{d}{dx}\frac{\partial f}{\partial y'} + \lambda\left(\frac{\partial g}{\partial y} - \frac{d}{dx}\frac{\partial g}{\partial y'}\right) = 0. \tag{1}
\end{equation}

Now consider the reciprocal problem: extremize $J$ subject to the constraint $I = \text{const}$. We introduce a Lagrange multiplier $\mu$ and form
\[
  K' = J + \mu I = \int_{x_1}^{x_2} [g + \mu\,f]\,dx.
\]
The Euler--Lagrange equation for $h' = g + \mu f$ is:
\begin{equation}
  \frac{\partial g}{\partial y} - \frac{d}{dx}\frac{\partial g}{\partial y'} + \mu\left(\frac{\partial f}{\partial y} - \frac{d}{dx}\frac{\partial f}{\partial y'}\right) = 0. \tag{2}
\end{equation}

Now, if $\lambda \neq 0$, we can divide equation (1) by $\lambda$:
\[
  \frac{1}{\lambda}\left(\frac{\partial f}{\partial y} - \frac{d}{dx}\frac{\partial f}{\partial y'}\right) + \frac{\partial g}{\partial y} - \frac{d}{dx}\frac{\partial g}{\partial y'} = 0.
\]
This is exactly equation (2) with $\mu = 1/\lambda$.

Therefore, any function $y(x)$ that satisfies the Euler--Lagrange equation for the original problem (extremize $I$ with $J$ fixed) also satisfies the Euler--Lagrange equation for the reciprocal problem (extremize $J$ with $I$ fixed), with the reciprocal Lagrange multiplier $\mu = 1/\lambda$. This establishes the reciprocity relationship.

\medskip
\noindent\textbf{Note:} The condition $\lambda \neq 0$ is essential. If $\lambda = 0$, the constraint $J = \text{const}$ plays no role, and $I$ is extremized unconditionally. In that case, the function that extremizes $I$ need not extremize $J$ subject to $I = \text{const}$.

\end{document}
